\documentclass[11pt,a4paper]{article}

\usepackage[margin=25mm]{geometry}
\usepackage[T1]{fontenc}
\usepackage[utf8]{inputenc}
\usepackage{lmodern}
\usepackage{microtype}
\usepackage{xcolor}
\usepackage{booktabs}
\usepackage{tabularx}
\usepackage{enumitem}
\usepackage[hidelinks]{hyperref}

\definecolor{aionblue}{HTML}{2368FF}
\definecolor{aiongreen}{HTML}{16834A}
\definecolor{aionred}{HTML}{C52F46}

\title{\textbf{AION Device Fabric}\\Games and Education Proof-of-Concept Supplement}
\author{Tessaris / AION Technical Record}
\date{29 August 2026\\Fabric implementation version 0.14.0}

\begin{document}
\maketitle

\begin{abstract}
This supplement records two application proofs implemented on top of the AION
Device Fabric: a governed cloud-gaming surface and an interactive Spanish
learning centre.  Neither proof requires additional television hardware.  The
existing LG webOS television is represented by a permissioned gateway, the
laptop operates as the mother node, and an iPhone provides a private control
surface.  The proofs demonstrate that a restricted television can become a
room-scale surface for composed intelligence, external services, local state,
and cross-device interaction.
\end{abstract}

\section{Common Fabric Substrate}

Both applications reuse the same Device Fabric rather than installing separate
systems.  The relevant topology is:

\begin{center}
\begin{tabular}{c}
User intent (voice, television remote, or private phone)\\
$\downarrow$\\
AION mother brain (planning, policy, memory, and audit)\\
$\downarrow$\\
Signed capability boundary and webOS gateway\\
$\downarrow$\\
LG television browser or third-party television application\\
$\updownarrow$\\
Private iPhone controller and, where supported, LG Mobile Gamepad
\end{tabular}
\end{center}

The mother brain retains the intelligence, policies, service state, and audit
evidence.  The television receives a temporary, token-protected presentation
surface.  The phone receives a separate private companion token and CSRF
secret.  Television display state therefore does not contain the webOS pairing
key, COMDEX credentials, payment data, or private service credentials.

\section{AION Games}

\subsection{Objective and demonstrated result}

AION Games composes an existing television, the NVIDIA GeForce NOW web client,
the AION webOS gateway, and a phone-based controller into a low-hardware cloud
gaming experience.  A command such as \emph{``Pilot, open games''} causes AION
to open the allowlisted GeForce NOW surface at
\texttt{https://play.geforcenow.com/}.  Game rendering and streaming remain an
NVIDIA service; AION supplies discovery, navigation, private text delivery,
controller composition, and governance.

The proof established that a compatible household television can expose a
large cloud-game catalogue without a console.  The user still requires a valid
NVIDIA account, any applicable subscription, a supported game entitlement,
sufficient network quality, and a controller path accepted by the television
and streaming service.

\subsection{Execution flow}

\begin{enumerate}[leftmargin=*,label=\arabic*.]
  \item A voice or phone command is parsed as the governed
        \texttt{open\_games} intent.
  \item The mother selects the enrolled LG gateway and verifies that the
        games capability is present.
  \item The webOS adapter opens the fixed GeForce NOW URL in the television
        browser and records a verified action receipt.
  \item The private phone companion supplies directional navigation, OK,
        back, pointer movement, pointer click, and game search controls.
  \item If typed sign-in is necessary, the user enters the username or
        password on the phone.  AION sends that value once to the currently
        focused television field.  The value is cleared from the form and is
        not placed in the Fabric audit ledger.
  \item CAPTCHA, two-factor authentication, legal consent, account selection,
        purchases, and the final sign-in decision remain explicit user acts.
  \item Once a stream begins, gameplay input can transfer to the LG Mobile
        Gamepad or another controller exposed by the platform.
\end{enumerate}

\subsection{Control and trust boundaries}

The GeForce NOW URL is allowlisted; arbitrary URL execution is not granted by
the games command.  The TV touchpad exists because some third-party pages
require pointer input, but it does not authorize AION to accept contracts,
purchase subscriptions, defeat access controls, or infer consent.  Account
secrets are intentionally excluded from normal Fabric state, telemetry,
screen projection, and audit receipts.

The phone controller is delivered on the private LAN through a rotating
session URL.  Each state-changing request carries a CSRF value, is rate
limited, returns to the mother brain, and passes through the same signed TV
policy used by voice commands.  A phone waking from sleep performs bounded
retries and a state refresh; loss of a transient browser connection therefore
does not create a new television enrollment.

\subsection{Implemented components}

\begin{table}[h]
\centering
\small
\begin{tabularx}{\textwidth}{@{}lX@{}}
\toprule
Component & Responsibility\\
\midrule
\texttt{runtime.py} & Governed companion commands, games capability routing,
and signed action policy.\\
\texttt{webos\_gateway.py} & LG pairing, application and browser launch,
remote-key delivery, pointer input, and verified receipts.\\
\texttt{companion.py} & Private remote, touchpad, one-shot credential delivery,
game search input, reconnect behaviour, and CSRF enforcement.\\
\texttt{tv\_canvas.py} & Room-scale AION Games status and phone-pairing
projection.\\
\bottomrule
\end{tabularx}
\end{table}

\subsection{Proof boundary and production work}

The proof demonstrates composition, not a replacement for NVIDIA, LG, or game
licensing.  Browser compatibility, codec support, controller exposure, network
latency, service-region restrictions, account status, and vendor UI changes
remain external dependencies.  A production release requires TLS on the phone
surface, durable device discovery, controller compatibility testing across LG
and Samsung models, accessible sign-in handoff, service health monitoring, and
clear subscription and data-use disclosures.

\section{AION Education: Spanish Learning Centre}

\subsection{Objective and demonstrated result}

The Spanish Learning Centre is a dedicated television application surface for
children approximately six to eight years old who speak English and are
learning Spanish.  It combines pictures, spoken Spanish, multiple-choice
interaction, immediate correction, gamified progress, adaptive repetition,
and short conversations.  The proof demonstrates that the same Fabric used to
control a restricted television can also host a stateful, personalised
learning experience.

\subsection{Learning model}

The current curriculum contains foundational greetings, kind words, food and
drink, home, school, people, and feelings.  Each word card contains a canonical
identifier, Spanish form, English meaning, category, and visual cue.  For
example, \emph{manzana} is presented together with an apple image cue.  Every
fourth round is a short conversation challenge rather than a vocabulary card.

Two anonymous local learner profiles are provided.  Each profile maintains:

\begin{itemize}[leftmargin=*]
  \item attempts and correct answers;
  \item stars and current answer streak;
  \item per-card mastery values;
  \item a 50-question session score; and
  \item the current lesson phase and round.
\end{itemize}

Cards with lower mastery are selected earlier.  A correct answer raises card
mastery and an incorrect answer reduces it slightly, causing weak material to
return sooner.  This is deliberately a small deterministic adaptive loop for
the proof; it does not claim to be a validated educational assessment model.

\subsection{Television interaction}

The dedicated application is served at the temporary Fabric path
\texttt{/tv/<session>/education}.  It exposes large focusable answer controls,
visual cues, score and question progress, and controller connection status.
Input is accepted from legacy webOS remote key codes, normal browser arrow and
Enter events, or the browser Gamepad API when LG exposes a compatible gamepad.

After selection, the interface:

\begin{enumerate}[leftmargin=*,label=\arabic*.]
  \item disables further answers for the current round;
  \item shows a green tick for a correct answer or a red cross for an
        incorrect answer;
  \item always identifies the correct answer;
  \item updates the score, stars, streak, attempts, and mastery;
  \item speaks the feedback; and
  \item advances automatically after speech finishes, with a bounded timeout
        fallback.
\end{enumerate}

Automatic advancement includes the expected round number.  The mother rejects
a duplicate or stale advance from a second browser or reconnecting controller,
preventing accidental skipped questions.  Periodic state refreshes also
preserve the currently focused answer instead of rebuilding focus every cycle.

\subsection{Natural speech pipeline}

The television does not need native speech synthesis.  It requests the current
bounded lesson phrase from the mother at
\texttt{/education/audio}.  The mother applies the following provider chain:

\begin{enumerate}[leftmargin=*]
  \item configured ElevenLabs \texttt{eleven\_multilingual\_v2} natural voice;
  \item configured OpenAI \texttt{gpt-4o-mini-tts} natural voice;
  \item local macOS Spanish speech as an offline fallback.
\end{enumerate}

Cloud requests contain only the fixed lesson phrase, limited to 240
characters.  They do not contain a child's name, recorded voice, microphone
audio, profile history, answer history, or household identity.  Returned audio
is limited to 2 MB and cached only in process memory with a bounded 96-entry
cache.  The television interface discloses that the natural learning voice is
AI-generated.  The audio control also permits an explicit \emph{Hear AI voice
again} action.

\subsection{State, actions, and security}

Persistent learning state is stored beneath the Fabric runtime directory in
\path{education/learning_centre.json}.  State writes use a temporary file
followed by an atomic replacement.  The public session projection removes the
hidden correct-answer index before it is sent to a client.

The television action endpoint accepts only \texttt{start}, \texttt{answer},
\texttt{next}, and \texttt{repeat}.  Requests must include the independent
\texttt{X-AION-Education} session secret and remain below the request-size
limit.  The phone uses the separate companion token and CSRF secret.  Education
actions are recorded in the Fabric audit chain without recording private
speech or credentials.

The child-safety posture currently enforces no child account, no advertising,
no open-web link, no purchase path, and no retained raw voice.  Learner labels
are abstract local profiles rather than legal identities.  The mother
microphone can remain physically closed while the lesson's output voice
continues to operate.

\subsection{Implemented components}

\begin{table}[h]
\centering
\small
\begin{tabularx}{\textwidth}{@{}lX@{}}
\toprule
Component & Responsibility\\
\midrule
\texttt{education.py} & Curriculum, anonymous profiles, mastery selection,
50-question sessions, scoring, feedback, and atomic persistence.\\
\texttt{tv\_canvas.py} & Dedicated visual learning application, remote and
gamepad mapping, protected actions, natural speech delivery, focus retention,
and automatic progression.\\
\texttt{companion.py} & Phone-based answers, lesson controls, private session
state, and wake/reconnect behaviour.\\
\texttt{runtime.py} & Education capability registration and governed routing
from voice or phone intent.\\
\texttt{cli.py} & Live service composition, television lesson handler, audit
events, and stale-advance protection.\\
\bottomrule
\end{tabularx}
\end{table}

\subsection{Verified proof status}

At the time of this record, the implementation has demonstrated:

\begin{itemize}[leftmargin=*]
  \item discovery, pairing, and verified control of the LG webOS television;
  \item launch of a dedicated learning page on the television;
  \item persistent independent progress for two anonymous learners;
  \item remote, phone, and compatible gamepad answer navigation;
  \item picture-supported vocabulary and conversation rounds;
  \item correct/incorrect visual feedback and explicit correction;
  \item automatic, duplicate-safe progression;
  \item natural Spanish speech with provider fallback;
  \item phone reconnection after browser sleep; and
  \item 127 passing automated Fabric tests with a valid audit chain.
\end{itemize}

\subsection{Remaining production work}

The next education phase should replace pictographic cues with a curated,
licensed illustration library; expand the curriculum and sentence model;
introduce age-appropriate difficulty bands; test pronunciation and regional
Spanish variants with educators; add parent-visible learning summaries;
measure retention across days; add accessibility controls; prefetch speech to
remove first-play latency; and conduct child-safety, privacy, and educational
efficacy review.  Natural speech providers also require explicit account,
quota, disclosure, retention, and offline-mode configuration in a consumer
installer.

\section{Conclusion}

The two proofs show the same architectural result from different directions.
AION Games composes a television, cloud service, phone, and controller into a
new entertainment surface.  AION Education composes the television, local
learner state, visual interaction, adaptive sequencing, and natural speech
into a new household learning surface.  The television is therefore not the
intelligence host by necessity; it is a governed room-scale endpoint of the
Fabric, while the mother brain supplies intelligence, memory, permissions,
recovery, and evidence.

\end{document}
